\documentclass[
  aps,prd,superscriptaddress,nofootinbib,preprint
]{revtex4-2}

\usepackage{amsmath,amssymb,amsthm}
\usepackage[colorlinks=true,linkcolor=blue,citecolor=blue,urlcolor=blue]{hyperref}

\newtheorem{theorem}{Theorem}
\newtheorem{lemma}{Lemma}
\newtheorem{proposition}{Proposition}
\newtheorem{corollary}{Corollary}
\theoremstyle{remark}
\newtheorem{remark}{Remark}

\newcommand{\CP}{\mathbb{CP}}
\newcommand{\cO}{\mathcal{O}}
\newcommand{\Spin}{\mathrm{Spin}}
\newcommand{\SU}{\mathrm{SU}}
\newcommand{\U}{\mathrm{U}}
\newcommand{\Tr}{\mathrm{Tr}}
\newcommand{\Res}{\mathrm{Res}}
\newcommand{\aut}{\mathfrak{aut}}
\newcommand{\sltwo}{\mathfrak{sl}(2,\mathbb C)}
\newcommand{\Ktr}{K_{\rm tr}}

\begin{document}

\title{Three fermion families from the automorphism algebra of their
own carrier, and the Killing-form Majorana contact}

\author{BoYu Chen}
\email{boypatrick.tw@gmail.com}
\affiliation{Independent Researcher, Hsinchu, Taiwan}

\date{\today}

\begin{abstract}
We prove three representation-theoretic statements about the family
structure of chiral matter.  (i)~Among the complex irreducible
representations of all simple compact Lie algebras, enumerated
exhaustively at dimension $\le 16$, the $\Spin(10)$ half-spinor is the
unique object whose restriction to the Standard-Model subalgebra
reproduces exactly one family; its sixteenth state --- a total
Standard-Model singlet, $\nu^c$ --- is thereby forced, not chosen.
(ii)~If the multiplicity space of families is required to carry, as its
covariance algebra, the automorphism algebra of its own geometric
carrier --- generated by the geometry rather than imposed upon it ---
then the number of families is bounded a priori by three, and the
requirement that the algebra admit a canonical nondegenerate invariant
bilinear form (the Cartan criterion) forces equality: the carrier is
the rational curve, the family space is its M\"obius algebra
$\sltwo$, and $N_{\rm fam}=\dim\sltwo=3$.  (iii)~The invariant Majorana
contact direction on the three-family space is unique up to scale and
equals the Killing form of that algebra,
$B=2\sqrt3\,\Ktr$ with $\Ktr^2=\tfrac13 I$; the same simplicity that
selects three families guarantees that the seesaw contact exists, is
unique, and is invertible.  No phenomenological claim is made: the
statements are unconditional consequences of the stated hypotheses, and
every step is machine-verified in exact arithmetic in a public
repository.
\end{abstract}

\maketitle

\section{Introduction}

Why the chiral matter of the Standard Model comes in three identical
families is an old structural question.  Three traditions address it.
The family-symmetry tradition imposes a horizontal group --- continuous
or discrete, $A_4$, $S_4$, $\SU(2)$, $\SU(3)$ and their relatives ---
on the multiplicity space and derives textures from its breaking
\cite{KingLuhn2013}; the group is chosen, not derived.  The
index-theorem tradition, descending from compactified string theory,
counts families as a topological index of an internal space
\cite{CandelasHorowitzStromingerWitten1985}; the number reflects the
choice of internal geometry.  A third, division-algebra tradition
relates the family triplication to exceptional algebraic structures
(octonionic trialities and their relatives; see
Ref.~\cite{BaezHuerta2010} for the algebraic setting of grand
unification); there the family symmetry arises as a discrete outer
automorphism of a division-algebra construction.

This Letter proves a different statement, which we have not found in
the literature in this combination.  We require the family multiplicity
space to be \emph{self-carried}: the symmetry acting on families is not
imposed and not inherited from a larger construction, but is the
automorphism algebra of the geometric object whose sections \emph{are}
the families.  Under this requirement the family number is bounded a
priori by three (Lemma~\ref{lem:ladder}); demanding in addition that
the family algebra admit a canonical nondegenerate invariant bilinear
form --- the algebraic ingredient a Majorana mass sector needs ---
selects, via the Cartan criterion, the unique case in which the bound
is saturated (Theorem~\ref{thm:threefam}); and the resulting contact
direction is then not merely \emph{like} a Killing form but \emph{is}
one (Theorem~\ref{thm:killing}).  Combined with an exhaustive
enumeration selecting the $\Spin(10)$ half-spinor as the unique
single-object realization of one family
(Theorem~\ref{thm:sixface}), the package answers ``why three'' and
``why a canonical seesaw direction'' with one piece of standard
mathematics used in a nonstandard place: simplicity of
$\mathfrak{sl}(2)$.

All results are elementary as mathematics; the claim is the assembly.
Every statement below is verified by standalone machine audits (exact
rational arithmetic for the enumeration; see
Sec.~\ref{sec:verification}), and nothing here asserts a
phenomenological fit: what the hypotheses do and do not deliver is
stated explicitly in Sec.~\ref{sec:notclaimed}.

\section{Hypotheses}
\label{sec:hypotheses}

Three hypotheses are used, stated here informally and used precisely
in the theorems.

\emph{(H1) Face projection with an observed-content anchor.}  One
family of chiral matter is realized as a single complex irreducible
representation $R$ of a simple compact group $G\supseteq G_{\rm SM}$,
whose restriction to $G_{\rm SM}=\SU(3)\times\SU(2)\times\U(1)_Y$
contains exactly the fifteen observed Weyl states of one family
(with their multiplicities), plus possibly Standard-Model singlets,
and no exotics; the gauge theory of $R$ must be anomaly-free.

\emph{(H2) Self-carried index.}  The family multiplicity space is
$V=H^0(\Sigma,\mathcal L)$ for a holomorphic line bundle
$\mathcal L$ on a compact Riemann surface $\Sigma$, and the covariance
algebra acting on families is required to be
$\aut(\Sigma)=H^0(\Sigma,T_\Sigma)$ acting on $V$ --- with the
self-carried clause $\mathcal L\simeq T_\Sigma$, so that the family
space is the automorphism algebra itself (the adjoint module), not a
module over an external symmetry.  An existence anchor $N_{\rm fam}\ge1$
is assumed.

\emph{(H3) Canonical contact.}  The family algebra admits a canonical
invariant symmetric bilinear form, nondegenerate on $V$ --- the
structure a family-covariant Majorana (seesaw) mass sector requires.

H1 is the standard grand-unification setting
\cite{Georgi1975,FritzschMinkowski1975} sharpened into an enumeration
hypothesis; H2 replaces an imposed family group by a generated one;
H3 is the minimal algebraic demand of a covariant seesaw.  None of the
three refers to a specific group, curve, or form: those are outputs.

\section{One family: the enumerated half-spinor}
\label{sec:sixteen}

\begin{lemma}[Chirality forces complexity]
\label{lem:complex}
If $\Res^G_{G_{\rm SM}}R$ satisfies the observed-content anchor, $R$ is
a complex representation of $G$.
\end{lemma}

\begin{proof}
Restriction commutes with conjugation, and the anchored content is not
self-conjugate: its hypercharge multiset contains $Y=1/6$ with
multiplicity six and $Y=-1/6$ with multiplicity zero, a mismatch that
extra self-conjugate singlets cannot repair.
\end{proof}

Since a maximal torus of $G_{\rm SM}$ embeds into one of $G$, the rank
of $G$ is at least four.  Only the $A_n$, $D_{2k+1}$, and $E_6$ series
admit complex irreducibles (the conjugation involution $-w_0$ is
computed from the Cartan matrix for every simple algebra of rank
$4$--$15$ and for $E_6,E_7,E_8,F_4$ in the machine audit; textbook
context in Ref.~\cite{Slansky1981}).  The Weyl dimension formula is
strictly increasing in each Dynkin label, so the following finite scan
is exhaustive.

\begin{proposition}[Enumeration at dimension $\le16$]
\label{prop:enumeration}
Among all simple compact groups of rank $\ge4$ (every algebra of rank
$4$--$15$ scanned directly; rank $\ge16$ contributes nothing, the
smallest rank-16 irreducibles having dimensions $17,33,32,32$), the
complex irreducible representations of dimension $15$ or $16$ are:
$\mathrm{Sym}^2 5$ of $\SU(5)$ and $\Lambda^2 6$ of $\SU(6)$
(both of which fail the content anchor: every embedding produces a
color sextet, octet, or a wrong antitriplet count); the fundamentals
of $\SU(15)$ and $\SU(16)$ (gauge-anomalous: on an exact integer
diagonal generator $\Tr T^3=2730$ and $3360$
\cite{GeorgiGlashow1972,Okubo1977}); and the $\Spin(10)$ half-spinor
$16$, together with conjugates.
\end{proposition}

\begin{theorem}[Enumerated one-family uniqueness]
\label{thm:sixface}
Under H1 the minimal-dimension realization exists, is unique, and is
the $\Spin(10)$ half-spinor, with
\[
  \Res^{\Spin(10)}_{G_{\rm SM}}16
  =Q\oplus L\oplus u^c\oplus d^c\oplus\nu^c\oplus e^c ,
\]
multiplicities $(6,2,3,3,1,1)$ as Weyl states.  The half-spinor is
anomaly-free: with the $16$ built explicitly from five fermionic
oscillators (chirality $=$ occupation parity),
$\Tr(\{T_a,T_b\}T_c)=0$ holds exactly over all $45^3$ generator
triples; the nonperturbative check is $\pi_4(\Spin(10))=0$
\cite{WittenSU2Anomaly1982}.
\end{theorem}

\begin{corollary}[The sixteenth state is a prediction]
\label{cor:nuc}
No fifteen-dimensional single-object realization of one family exists;
the sixteenth state --- a total Standard-Model singlet, $\nu^c$ --- is
forced.  The essential seesaw ingredient is an output of H1, not an
input.  (Instructive near-misses: $\mathfrak{so}(9)$ has a \emph{real}
spinor of dimension exactly $16$ and $\mathfrak{so}(15)$ a real vector
of dimension $15$; both fail Lemma~\ref{lem:complex}.  Chirality, not
size, selects $\Spin(10)$.)
\end{corollary}

On the surviving object $Y=T_{3R}+(B-L)/2$, and the weight-level
traces give $\Tr Y^2=10/3$, $\Tr T_{3L}^2=2$, hence the canonical
normalization $k_Y=5/3$.

\section{Three families: the self-carried bound and the Cartan
selection}
\label{sec:three}

\begin{lemma}[Genus ladder]
\label{lem:ladder}
For a compact Riemann surface $\Sigma_g$,
\[
  h^0(\Sigma_g,T_{\Sigma_g})=
  \begin{cases} 3, & g=0,\\ 1, & g=1,\\ 0, & g\ge2. \end{cases}
\]
Consequently, under H2, $N_{\rm fam}\le3$ \emph{a priori}, with
equality if and only if the family curve is rational.
\end{lemma}

\begin{proof}
$T_{\Sigma_g}=K^{-1}_{\Sigma_g}$ has degree $2-2g$.  For $g=0$,
Riemann--Roch gives $h^0-h^1=3$ and Serre duality gives
$h^1=h^0(K^2)=h^0(\cO(-4))=0$; for $g=1$ the tangent bundle is trivial;
for $g\ge2$ its degree is negative
\cite{GriffithsHarris1978,LawsonMichelsohn1989}.
\end{proof}

\begin{lemma}[Cartan selection]
\label{lem:cartan}
Among the cases $g\in\{0,1\}$ allowed by the existence anchor, only
$g=0$ satisfies H3.
\end{lemma}

\begin{proof}
For $g=1$ the automorphism algebra is one-dimensional abelian; the
Killing form of an abelian algebra vanishes identically, so by the
Cartan criterion no canonical nondegenerate invariant form exists.
For $g=0$, $\aut(\CP^1)=\sltwo$ is simple and its Killing form is
nondegenerate.
\end{proof}

\begin{theorem}[Three families]
\label{thm:threefam}
Under H2 and H3: the family curve is $\Sigma\simeq\CP^1$; the family
space is canonically the M\"obius algebra
$V=H^0(\CP^1,T_{\CP^1})\simeq\sltwo$; the carrier is
$T_{\CP^1}\simeq\cO(2)$; and $N_{\rm fam}=\dim\sltwo=3$.
\end{theorem}

Two structural remarks.  First, the zero divisor of any regular
semisimple element $X_h\in V$ (any M\"obius flow) is a sum of two
distinct points with
$\deg\mathrm{div}(X_h)=c_1(T)[\CP^1]=\chi(S^2)=2$: the ``two centers''
often postulated in family-geometry constructions are the two fixed
points of a Cartan flow --- Poincar\'e--Hopf, not an assumption.
Second, the self-carried clause in H2 does real selective work: for
$g=0$ one has $h^0(K^{-m})=2m+1$, so ``some power of the anticanonical
bundle'' would allow any odd family number.  Only for $m=1$ is the
family space itself the algebra generating the family covariance ---
the adjoint module --- rather than a module over an external symmetry.
This is precisely what separates the present mechanism from the
imposed-symmetry tradition.

\section{The contact is the Killing form}
\label{sec:killing}

A family-covariant Majorana sector requires an invariant symmetric
bilinear form on $V$.

\begin{lemma}[Uniqueness up to scale]
\label{lem:unique}
$\mathrm{Sym}^2(\mathrm{Sym}^2\mathbb C^2)\simeq
\mathrm{Sym}^4\mathbb C^2\oplus\mathrm{Sym}^0\mathbb C^2$, so the
invariant symmetric bilinear form on the adjoint (spin-1) module of
$\sltwo$ is unique up to scale \cite{FultonHarris1991}.
\end{lemma}

\begin{theorem}[Killing contact]
\label{thm:killing}
In the spherical orthonormal family basis $(T_{+1},T_0,T_{-1})$,
$T_{\pm1}=\mp J_\pm/\sqrt2$, $T_0=J_z$ (spin-1 matrices), the Killing
form of the family algebra is
\[
  B=\begin{pmatrix}0&0&-2\\0&2&0\\-2&0&0\end{pmatrix}
  =2\sqrt3\,\Ktr,\qquad
  \Ktr=\frac1{\sqrt3}
  \begin{pmatrix}0&0&-1\\0&1&0\\-1&0&0\end{pmatrix},
\]
with $\Ktr^2=\tfrac13 I$ and $\Ktr^{-1}=3\Ktr$.  It exists and is
nondegenerate because $\sltwo$ is simple (Cartan criterion), and it is
the unique invariant contact direction by Lemma~\ref{lem:unique}.
\end{theorem}

\begin{proof}
The adjoint representation of $\sltwo$ is the spin-1 representation,
so $B(x,y)=\Tr_{3\times3}(R_1(x)R_1(y))$ exactly.  With
$\Tr J_z^2=2$, $\Tr J_+J_-=4$, $\Tr J_\pm^2=0$:
$B(T_0,T_0)=2$, $B(T_{+1},T_{-1})=-\tfrac12\Tr J_+J_-=-2$,
$B(T_{\pm1},T_{\pm1})=0$.
\end{proof}

\begin{remark}[One principle, two consequences]
The same simplicity that selects $g=0$ --- and hence exactly three
families --- guarantees that the seesaw contact direction exists, is
unique, and is invertible.  In a hypothetical $g=1$ branch the Majorana
contact would have no canonical direction at all.  ``Exactly three
families'' and ``a working invariant seesaw contact'' are two faces of
one algebraic fact.
\end{remark}

\section{What is not claimed}
\label{sec:notclaimed}

The theorems are unconditional consequences of H1--H3; the hypotheses
themselves are physical choices, and we do not claim them forced.  In
particular: no claim is made that a specific grand-unified dynamics
realizes the package, no flavor fit or proton-decay phenomenology is
asserted here, and the \emph{coefficient} multiplying the Killing
contact in any realized Majorana sector is an input --- in the
companion framework this last statement is itself a theorem (the
coefficient is provably underivable from hypotheses of the above
kind), and the accompanying phenomenological audits are reported
elsewhere.  The content of this Letter is the rigid part: one family
forces the half-spinor and its sixteenth state; self-carried covariance
bounds the family number by three; the Cartan criterion saturates the
bound and hands the seesaw its unique canonical direction.

\section{Machine verification}
\label{sec:verification}

Every claim above is verified by standalone audits (plain
\texttt{python3}/\texttt{numpy}; the enumeration in exact rational
arithmetic): the $-w_0$ involutions for all simple algebras of rank
$4$--$15$ and $E_6,E_7,E_8,F_4$; the full dimension-$\le16$ scan with
the $\mathfrak{so}(9)$/$\mathfrak{so}(15)$ near-misses; the exact
$\SU(15)/\SU(16)$ cubic traces; the exactly vanishing cubic trace
tensor of the oscillator-built $16$ over all $45^3$ triples and its
demicube weight system; the hypercharge traces and $k_Y=5/3$; the
one-dimensionality of the invariant-contact solution space; and the
identity $B=2\sqrt3\,\Ktr$.  The audit scripts and machine-readable
ledgers (68 checks) are available in the public repository
\href{https://github.com/boypatrick/GUT_routeA}{github.com/boypatrick/GUT\_routeA}.

\begin{acknowledgments}
The machine audits were developed with AI assistance; all mathematical
claims were verified by the scripts cited above.
\end{acknowledgments}

\bibliography{../../paper/refs}

\end{document}
